\section{PART I --- INFORMATION THEORY}

\subsection{1. Entropy and Mutual Information}
\begin{itemize}
\item Shannon entropy; joint, conditional, mutual information
\item Relative entropy (KL divergence); properties; chain rules
\item Differential entropy; Gaussian channels; entropy of mixtures
\item Data processing inequality; Fano's inequality
\end{itemize}

\subsection{2. Lossless Source Coding}
\begin{itemize}
\item Asymptotic equipartition property (AEP); typical sequences
\item Shannon's source coding theorem; prefix codes; Huffman coding
\item Arithmetic coding; Lempel-Ziv universal compression
\item Kolmogorov complexity (brief); connections to learning
\end{itemize}

\subsection{3. Channel Coding}
\begin{itemize}
\item Channel capacity; Shannon's noisy-channel theorem; proof via typicality
\item Gaussian channel; water-filling; multi-antenna (MIMO) brief
\item Linear codes; Reed-Solomon (from finite fields, Course 2); LDPC
\item Polar codes and capacity-achieving constructions (brief)
\end{itemize}

\subsection{4. Information-Theoretic Inequalities}
\begin{itemize}
\item Log-sum inequality; convexity of KL divergence
\item Fisher information; Cram\'er-Rao bound; van Trees inequality
\item Entropy power inequality; de Bruijn's identity
\item Connections to Course 1 (Sections 4, 28) and Course 5
\end{itemize}

\subsection{5. Information Theory and Statistics}
\begin{itemize}
\item Hypothesis testing; Neyman-Pearson lemma; error exponents
\item Chernoff bound; Stein's lemma; Hoeffding bound
\item Method of types; empirical entropy; connection to large deviations
\item Universal hypothesis testing; I-projection
\end{itemize}

\subsection{6. Information Theory and Machine Learning}
\begin{itemize}
\item Information bottleneck; rate-distortion perspective on representations
\item MI estimation with neural networks (MINE, NWJ bounds)
\item Differential privacy: formal definition; mechanisms; composition
\item Rate-distortion theory and representation learning; connections to VAEs
\end{itemize}

\section{PART II --- STATISTICAL ESTIMATION THEORY}

\subsection{7. Classical Estimation}
\begin{itemize}
\item Maximum likelihood; consistency; asymptotic normality; efficiency
\item Cram\'er-Rao bound; Fisher information matrix; Bhattacharyya bound
\item M-estimators; robustness; breakdown point; influence function
\item Bayesian estimation; conjugate priors; posterior; minimax Bayes
\end{itemize}

\subsection{8. Exponential Families and Sufficient Statistics}
\begin{itemize}
\item Natural parameters; log-partition function; moment parameterization
\item Sufficient statistics; Fisher-Neyman factorization; completeness
\item Rao-Blackwell theorem; Lehmann-Scheff\'e theorem
\item Connections to information geometry (Course 7, Section 6)
\end{itemize}

\subsection{9. Minimax Theory and Optimal Rates}
\begin{itemize}
\item Minimax risk; parameter space complexity; Le Cam's two-point method
\item Assouad's lemma; Fano's inequality applied to estimation
\item Minimax rates for nonparametric regression and density estimation
\item Adaptive estimation: Lepski's method; oracle inequalities; wavelet thresholding
\end{itemize}

\section{PART III --- STATISTICAL LEARNING THEORY}

\subsection{10. PAC Learning Framework}
\begin{itemize}
\item PAC model; sample complexity; $\varepsilon$-$\delta$ learnability; realizable and agnostic
\item VC dimension; growth function; Sauer-Shelah lemma
\item Fundamental theorem of statistical learning: VC characterizes PAC learnability
\item Examples: halfspaces, decision trees, neural networks
\end{itemize}

\subsection{11. Rademacher Complexity and Uniform Convergence}
\begin{itemize}
\item Rademacher complexity; symmetrization lemma
\item Covering numbers; metric entropy; Dudley's chaining integral
\item Generalization bounds via Rademacher complexity; data-dependent bounds
\item PAC-Bayes bounds; KL-regularized objectives; connections to Bayesian methods
\end{itemize}

\subsection{12. Stability, Compression, and Generalization}
\begin{itemize}
\item Algorithmic stability: uniform stability; connection to generalization
\item On-average stability; SGD is stable; stability of noisy algorithms
\item Compression-based bounds; mutual information bounds
\item Connections to regularization; data augmentation as stability
\end{itemize}

\subsection{13. Kernel Methods and RKHS}
\begin{itemize}
\item Reproducing kernel Hilbert spaces; Mercer's theorem; feature maps
\item Kernel ridge regression; SVMs; representer theorem
\item Spectral analysis of kernels; eigenvalue decay; approximation rates
\item Random features; Nystr\"om approximation; scalable kernel methods
\end{itemize}

\subsection{14. High-Dimensional Statistics}
\begin{itemize}
\item Curse of dimensionality; concentration of measure; geometric results
\item Lasso: restricted eigenvalue condition; oracle inequalities; de-biasing
\item Covariance estimation; sample covariance in high dimensions; shrinkage
\item Connections to compressed sensing (Course 1, Section 33)
\end{itemize}

\subsection{15. Neural Network Theory}
\begin{itemize}
\item Overparameterization; benign overfitting; interpolation threshold
\item Neural tangent kernel (NTK): linearization; kernel regime; lazy training
\item Mean field theory of two-layer networks; Wasserstein gradient flow
\item Double descent: risk curve; model and epoch double descent
\item Implicit regularization of gradient descent; connections to sparsity
\end{itemize}

\subsection{16. Online Learning and Bandits}
\begin{itemize}
\item Online convex optimization; regret; minimax regret; expert advice
\item Hedge (Multiplicative Weights) algorithm; regret bounds; connections to entropy
\item Multi-armed bandits: UCB; Thompson sampling; lower bounds; Lai-Robbins
\item Contextual bandits; linear bandits; connections to RL
\end{itemize}
